与余积相容的分次：III，§ 11，第 1 号。

与代数结构相容的分次：III，§ 3，第 1 号。

诱导分次，商分次：II，§ 11，第 3 号。

\(\Delta\) 型分次：II，§ 11，第 1 号。

分次，偏分次，全分次：II，§ 11，第 1 和第 2 号。

平凡分次：II，§ 11，第 1 号。

格拉斯曼簇：III，§ 11，第 13 号。

格拉斯曼关系：III，§ 11，第 13 号。

两个整数的最大公因数（g.c.d.）：I，§ 8，第 11 号。

群：I，§ 2，第 3 号。

环的加法群：I，§ 8，第 1 号。

仿射群：I，§ 9，第 4 号。

交错群：I，§ 5，第 7 号。

双循环群：I，§ 6，习题 26。

双分次群：II，§ 11，第 2 号。

两个同态的重合子群：I，§ 4，第 8 号。

换位子群：I，§ 6，第 2 号。

循环群：I，§ 4，第 10 号。

由生成元和关系定义的群：I，§ 7，第 6 号。

导群：I，§ 6，第 2 号。

二面体群：I，§ 6，习题 4。

有限群，无限群：I，§ 4，第 1 号。

有限生成群：I，§ 7，习题 16。

有限表现群：I，§ 7，习题 16。

集合上的自由交换群：I，§ 7，第 7 号及 II，§ 1，第 11 号。

集合上的自由群：I，§ 7，第 5 号。

分次群：II，§ 11，第 1 号。

线性群：II，§ 2，第 6 号。

极小单群：I，§ 6，习题 27。

单生成群：I，§ 4，第 10 号。

环的乘法群：I，§ 8，第 1 号。

幂零群，n 类幂零群：I，§ 6，第 3 号。

差群，分式群：I，§ 2，第 4 号。

指数型群：I，§ 7，习题 39。

忠实作用的群：I，§ 5，第 1 号。

自由作用的群：I，§ 5，第 4 号。

单传递作用的群：I，§ 5，第 6 号。

传递作用的群：I，§ 5，第 5 号。

p-群：I，§ 6，第 5 号。

射影群：II，§ 9，第 10 号。

剩余有限群：I，§ 5，习题 5。

可解群，类 n 的可解群：I，§ 6，第 4 号。

特殊线性群：III，§ 8，第 9 号。

超可解群：I，§ 6，习题 26。

对称群：I，§ 4，第 1 号。

幺模群：III，§ 8，第 9 号。

带算子群，阿贝尔的，交换的：I，§ 4，第 2 号。

带算子群：I，§ 4，第 2 号。

带算子群的积：I，§ 4，第 8 号。

带算子群的积（或内积），一族商群的积：I，§ 4，第 8 号。

带算子群的商：I，§ 4，第 4 号。

带算子单群：I，§ 4，第 4 号。

群胚：I，§ 4，习题 23。

霍尔定理：I，§ 6，习题 10。

哈默尔基：II，§ 7，第 1 号。

哈密顿四元数：III，§ 2，第 5 号。

分次群中的齐次元素：II，§ 11，第 1 号。

齐次 G-集：I，§ 5，第 5 号。

齐次线性方程，线性方程组：II，§ 2，第 8 号。

度为 \(p\) 的形式幂级数中的齐次子集：III，§ 2，第 11 号。

度为 \(p\) 形式幂级数中关于某些未定元的齐次子集：III，§ 2，第 11 号。

同态，代数：III，§ 1，第 1 号。

同态，A-模，A-同态：II，§ 1，第 2 号。

同态，中心环：II，§ 5，第 3 号。

同态，交叉：I，§ 6，习题 7。

同态，本性的：II，§ 2，习题 15。

两个合成律的同态：I，§ 1，第 1 号。

同态，分次：II，§ 11，第 2 号。

同态，分次代数：III，§ 3，第 1 号。

同态，群：I，§ 4，第 1 号。

同态，幺半群：I，§ 2，第 1 号。

同态，M-集：I，§ 5，第 1 号。

同态，多重模：II，§ 1，第 14 号。

带算子群的同态：I，§ 4，第 2 号。

\(\phi\) -同态：I，§ 3，第 1 号。

同态，投影：I，§ 4，第 8 号。

同态，环：I，§ 8，第 4 号。

同态，平凡：I，§ 2，第 1 号。

同态，保单位的：I，§ 2，第 1 号。

同态，保单位代数：III，§ 1，第 1 号。

位似：I，§ 4，第 2 号和 II，§ 1，第 1 号。

位似，中心：II，§ 1，第 2 号和 III，§ 9，习题 6。

超平面，仿射：II，§ 9，第 3 号。

无穷远超平面：II，§ 9，第 8 号。

向量空间中过 0 的超平面：II，§ 7，第 3 号。

超平面，射影：II，§ 9，第 7 号。

超平面，射影，取作无穷远超平面：II，§ 9，第 10 号。

理想，判别式：III，§ 9，第 5 号。

理想，分次：II，§ 11，第 3 号和 III，§ 3，第 2 号。

理想，不可约双边：I，§ 8，习题 11。

理想，左理想、右理想、双边理想（代数中）：III，§ 1，第 2 号。

理想，左理想、右理想、双边理想（环中）：I，§ 8，第 6 号。

理想，极大：I，§ 8，第 6 号。

理想，素：I，§ 9，第 3 号。

理想，主：I，§ 8，第 6 号。

关系子理想：III，§ 2，第 8 号。

理想，零：I，§ 8，第 6 号。

幂等元：I，§ 1，第 4 号。

恒等式，雅可比，行列式的子式之间的：III，§ 11，习题 9。

恒等式，多项式：III，§ 2，第 9 号。

单位元：I，§ 2，第 1 号。

恒等式，雅可比：III，§ 10，第 6 号。

恒等式，Redei：III，§ 8，习题 25。

射影线性簇在射影映射下的像（原像）：II，§ 9，第 10 号。

同态的像：II，§ 1，第 3 号。

未定元，指标为 \(i\) 的未定元：I，§ 7，第 5 号和 III，§ 2，第 7 号。

子群的指数：I，§ 4，第 4 号。

诱导作用：I，§ 3，第 2 号。

诱导分次：II，§ 11，第 3 号。

诱导 K'-结构：II，§ 8，第 1 号。

诱导律：I，§ 1，第 4 号。

无穷远，超平面：II，§ 9，第 8 号。

无穷远，点：II，§ 9，第 8 号。

群的内自同构：I，§ 5，第 3 号。

环的内自同构：I，§ 8，第 4 号。

内乘，左、右：III，§ 11，第 6 和第 7 号。

整数，负、正、严格负、严格正：I，§ 2，第 5 号。

整数，有理：I，§ 2，第 5 号。

有理整数模 \(a\) （环）：I，§ 8，第 11 号。

整数，互素：I，§ 8，第 11 号。

整环：I，§ 9，第 2 号。

整性，整区：I，§ 9，第 2 号。

内直积，限制和：I，§ 4，第 9 号。

不变元：I，§ 5，第 2 号。

不变子群，稳定子群：I，§ 4，第 4 号。

逆元，左逆元，右逆元，元素：I，§ 2，第 3 号。

线性映射的逆，左逆，右逆：II，§ 1，第 9 号。

逆向极限：见极限，逆向。

逆向系统：见系统，逆向。

双代数中的反演：III，§ 11，习题 4。

置换的逆序：I，§ 5，第 7 号。

可逆，左可逆，右可逆，元素：I，§ 2，第 3 号。

可逆，左可逆，右可逆，线性映射：II，§ 1，第 9 号。

可逆方阵：II，§ 10，第 7 号。

不可约理想：I，§ 8，习题 11。

等权元素：II，§ 11，第 1 号。

同构的广群：I，§ 1，第 1 号。

同构，广群：I，§ 1，第 1 号。

雅可比恒等式：III，§ 10，第 6 号。

Jordan-Hölder 列：I，§ 4，第 7 号。

Jordan-Hölder 定理：I，§ 4，第 7 号。

并置：I，§ 7，第 2 号。

群同态的核：I，§ 4，第 5 号。

线性映射的核：II，§ 1，第 3 号。

克罗内克符号：II，§ 1，第 11 号。

克鲁尔定理：I，§ 8，第 6 号。

拉普拉斯展开：III，§ 8，第 6 号。

律与等价关系，相容：I，§ 1，第 6 号。

律，结合：I，§ 1，第 3 号。

律，交换：I，§ 1，第 5 号。

律，群：I，§ 4，第 1 号。

律，诱导：I，§ 1，第 4 号。

律，左分配，右分配，关于另一个律分配：I，§ 3，第 4 和 5 号。

并非处处定义的律：I，§ 1，第 1 号。

合成律：I，§ 1，第 1 号。

幺半群在集合上的左、右作用律：I，§ 5，第 1 号。

与作用相伴的右、左作用律：I，§ 3，第 1 号。

律，相反：I，§ 1，第 1 号。

律，商：I，§ 1，第 6 号。

加法、乘法记法的律：I，§ 1，第 1 号。

两个整数的最小公倍数（l.c.m.）：I，§ 8，第 11 号。

Leibniz 公式：III，§ 10，第 4 号。

引理，扎森豪斯的：I，§ 4，第 7 号。

在融合和中分解的长度：I，§ 7，第 3 号。

群的长度：I，§ 4，第 7 号。

模的长度：II，§ 1，第 10 号。

自由群中元素的长度：I，§ 7，习题 19。

自由广群中元素的长度：I，§ 7，第 1 号。

字的长度：I，§ 7，第 2 号。

极限，正向，\(\mathbf{A}_{\alpha}\) -代数的：III，§ 1，第 6 号。

极限，正向，\(\mathbf{A}_{\alpha}\) -模的：II，§ 6，第 2 号。

极限，正向，作用的：I，§ 10，第 4 号。

极限，正向，分次代数的：III，§ 3，第 3 号。

极限，正向，群、带算子群、幺半群的：I，§ 10，第 3 号。

极限，正向，广群的：I，§ 10，第 3 号。

极限，正向，环的：I，§ 10，第 3 号。

极限，逆向，\(\mathbf{A}_{\alpha}\) -代数的：III，§ 1，第 6 号。

极限，逆向，\(\mathbf{A}_{\alpha}\) -模的：II，§ 6，第 1 号。

极限，逆向，群、带算子群、幺半群的：I，§ 10，第 1 号。

极限，逆向，广群的：I，§ 10，第 1 号。

极限，逆向，环的：I，§ 10，第 1 号。

直线，仿射：II，§ 9，第 1 和 3 号。

向量空间中的直线（过 0）：II，§ 7，第 3 号。

直线，射影：II，§ 9，第 5 号。

线性方程，方程组：II，§ 2，第 8 号。

线性方程，方程组，齐次：II，§ 2，第 8 号。

线性形式：II，§ 2，第 3 号。

线性群：II，§ 2，第 6 号。

线性群，特殊的：III，§ 8，第 9 号。

线性映射：II，§ 1，第 2 号。

线性映射，函数，仿射的：II，§ 9，第 4 号。

线性映射，射影的：II，§ 9，第 10 号。

线性问题：II，§ 2，第 8 号。

线性簇：II，§ 9，第 3、7 和 11 号。

线性簇，仿射的：II，§ 9，第 3 号。

线性簇，射影的：II，§ 9，第 7 和 11 号。

线性相关，线性无关，元素：II，§ 1，第 11 号。

线性不相交子代数：III，§ 4，第 4 号。

线性性，延拓原理：II，§ 1，第 7 号。

广群：I，§ 1，第 1 号。

广群，结合的：I，§ 1，第 3 号。

广群，交换的：I，§ 1，第 5 号。

由生成元和关系定义的广群：I，§ 7，第 1 号。

广群，自由的，在一个集合上的：I，§ 7，第 1 号。

到一个广群中的映射所成的广群：I，§ 1，第 1 号。

广群，相反的：I，§ 1，第 1 号。

广群，积：I，§ 1，第 1 号。

广群，商：I，§ 1，第 6 号。

广群，含幺的：I，§ 2，第 1 号。

广群，同构的：I，§ 1，第 1 号。

映射，仿射线性的，仿射的：II，§ 9，第 4 号。

映射，交错多重线性的：III，§ 7，第 4 号。

映射，双加法的，Z-双线性的：II，§ 3，第 1 号。

映射，C-双线性的：II，§ 3，第 5 号。

与一个作用相容的映射：I，§ 3，第 1 号。

与一个幺半群的作用相容的映射：I，§ 5，第 1 号。

关于一个变量分配的映射：I，§ 3，第 4 号。

映射，线性的，A-线性的：II，§ 1，第 2 号。

映射，线性的， \(\mathrm{A_s(T)} \to \mathrm{E}\) ，由映射 \(\mathrm{T} \to \mathrm{E}\) 确定：II，§ 1，第 11 号。

映射，线性的，与一个仿射映射相伴的：II，§ 9，第 4 号。

映射，多重加法的，Z-多重线性的：II，§ 3，第 9 号。

映射，C-多重线性的：II，§ 3，第 9 号。

映射，轨道的：I，§ 5，第 4 号。

映射，射影线性的，射影的：II，§ 9，第 10 号。

映射，右可逆线性的，左可逆线性的：II，§ 1，第 9 号。

映射，半线性的：II，§ 1，第 13 号。

映射，对称多重线性的：III，§ 6，第 3 号。

仅列、行顺序不同的矩阵：II，§ 10，第 9 号。

矩阵，等价的：II，§ 10，第 9 号。

矩阵，相似的方阵：II，§ 10，第 9 号。

矩阵，一个可逆矩阵的逆步矩阵：II，§ 10，第 7 号。

矩阵，对角的：II，§ 10，第 7 号。

矩阵，空的：II，§ 10，第 1 号。

矩阵，可逆的：II，§ 10，第 7 号。

矩阵，下三角的，上三角的：II，§ 10，第 7 号。

矩阵，类型为 \((p,q)\) 的矩阵：II，§ 10，第 1 号。

矩阵，单项的：II，§ 10，第 7 号。

通过给一个矩阵加边得到的矩阵：II，§ 10，第 1 号。

通过删去列、行得到的矩阵：II，§ 10，第 1 号。

一个线性映射关于两个基的矩阵：II，§10，第 4 号。

一个线性方程组的矩阵：II，§ 10，第 4 号。

一个元素关于一个基的矩阵：II，§ 10，第 4 号。

一个自同态关于一个基的矩阵：II，§ 10，第 7 号。

一个置换的矩阵：II，§ 10，第 7 号。

一个半线性映射关于两个基的矩阵：II，§10，第 6 号。

从一个基到另一个基的过渡矩阵：II，§ 10，第 8 号。

矩阵，纯量的：II，§ 10，第 7 号。

矩阵，方阵，\(n\) 阶方阵：II，§ 10，第 7 号。

矩阵，幺模：III，§ 8，第 4 号。

矩阵单位：II，§ 10，第 3 号。

对角线以下（以上）全为零的矩阵：II，§ 10，第 7 号。

矩阵，零：II，§ 10，第 2 号。

极大理想：I，§ 8，第 6 号。

G-均值：I，§ 6，习题 8。

极小正规稳定子群：I，§ 4，习题 15。

极小单群：I，§ 6，习题 27。

子式，矩阵的 \(p\) 阶子式：III，§ 8，第 5 号。

子式，互补：III，§ 8，第 6 号。

混合张量：III，§ 5，第 6 号。

模，双分次：II，§ 11，第 2 号。

模，可除：II，§ 7，习题 33。

模，对偶：II，§ 2，第 3 号。

模，忠实：II，§ 1，第 12 号。

模，忠实的，与一个模相伴的：II，§ 1，第 12 号。

模，自由：II，§ 1，第 11 号。

模，分次自由：II，§ 11，第 2 号。

模，分次，具有正次数的分次模：II，§ 11，第 2 号。

模，分次商：II，§ 11，第 3 号。

模，不可分解：II，§ 2，习题 21。

模，内射：II，§ 2，习题 11。

模，左模、右模、A-模：II，§ 1，第 1 号。

模，单生成的：II，§ 1，第 12 号。

有限长度模：II，§ 1，第 10 号。

形式线性组合模：II，§ 1，第 11 号。

线性关系模：II，§ 1，第 11 号。

代数上的模：III，§ 4，第 3 号。

模，积：II，§ 1，第 5 号。

模，投射：II，§ 2，第 2 号。

模，商：II，§ 1，第 3 号。

模，自反：II，§ 2，第 7 号。

模，无挠，整环上的：II，§ 7，第 10 号。

模，挠模，整环上的：II，§ 7，第 10 号。

单生成群：I，§ 4，第 10 号。

单生成模：II，§ 1，第 2 号。

幺半群：I，§ 2，第 1 号。

由生成元和关系定义的幺半群：I，§ 7，第 2 号。

幺半群，自由，集合上的：I，§ 7，第 2 号。

差幺半群：I，§ 2，第 4 号。

分母在 S 中的分式幺半群：I，§ 2，第 4 号。

在一个集合上忠实作用的幺半群：I，§ 5，第 1 号。

幺半和：I，§ 7，第 3 号。

单项式：III，§ 2，第 9 号。

态射，代数：III，§ 1，第 1 号。

态射，A-模：II，§ 1，第 2 号。

态射，余代数：III，§ 11，第 1 号。

态射，分次双代数、左分次双代数态射：III，§ 11，第 4 号。

态射，分次余代数：III，§ 11，第 1 号。

态射，广群：I，§ 1，第 1 号。

态射，幺半群：I，§ 2，第 1 号。

扩张的态射：I，§ 6，第 1 号。

态射，环：I，§ 8，第 4 号。

态射，保单位代数：III，§ 4，第 5 号。

态射，含幺广群：I，§ 2，第 1 号。

\(\phi\) -态射（ \(\phi\) 为一个算子幺半群到另一个的同态）：I，§ 5，第 1 号。

\(\phi\) -态射（ \(\phi\) 为一个算子集合到另一个的映射）：I，§ 3，第 1 号。

M-态射（M 为一个算子幺半群）：I，§ 5，第 1 号。

Ω-态射（Ω 为一个算子集合）：I，§ 3，第 1 号。

多重加法、Z-多重线性映射：II，§ 3，第 9 号。

幺半群的自由代数中的多重次数：III，§ 2，习题 13。

重指标：I，§ 7，第 7 号。

多重线性形式：II，§ 3，第 9 号。

C-多重线性映射：II，§ 3，第 9 号。

多重模：II，§ 1，第 14 号。

多重模，商：II，§ 1，第 14 号。

倍元，左、右：I，§ 8，第 1 号。

乘法：I，§ 1，第 1 号。

代数中的乘法：III，§ 1，第 1 号。

乘法表：III，§ 1，第 7 号。

环的乘法群：I，§ 8，第 1 号。

元素的负元：I，§ 2，第 3 号。

负有理整数：I，§ 2，第 5 号。

负有理数：I，§ 9，第 4 号。

S-邻近元素：I，§ 7，习题 18。

尼尔森-施赖埃尔定理：I，§ 7，习题 20。

幂零群，类为 \(n\) 的幂零群：I，§ 6，第 3 号。

范数，凯莱：III，§ 2，第 4 号。

二次代数中的范数：III，§ 2，第 3 号。

K-代数中元素关于 K 的范数：III，§ 9，第 3 号。

标量关于模的范数：III，§ 9，第 1 号。

正规稳定子群：I，§ 4，第 4 号。

正规子群：I，§ 4，第 4 号。

正规化子：I，§ 5，第 3 号。

使子集正规化（元素，子集）：I，§ 5，第 3 号。

零元素：I，§ 2，第 1 号。

数，素数：I，§ 4，第 10 号。

数，有理数：I，§ 9，第 4 号。

（有理）数，负的，正的，严格负的，严格正的：I，§ 9，第 4 号。

分子：I，§ 2，第 4 号。

八元数，凯莱：III，附录，第 3 号。

型 \((\alpha, \beta, \gamma, \delta)\) 的八元数（代数）：III，附录，第 3 号。

奇置换：I，§ 5，第 7 号。

通过左、右平移的作用：I，§ 5，第 1 号。

作用，左、右（律）：I，§ 5，第 1 号。

幺半群的左、右作用：I，§ 5，第 1 号。

作用，单传递的：I，§ 5，第 6 号。

作用，传递的：I，§ 5，第 5 号。

作用，平凡的：I，§ 5，第 2 号。

算子：I，§ 3，第 1 号。

反代数：III，§ 1，第 1 号。

反余代数：III，§ 11，第 1 号。

反合成律：I，§ 1，第 1 号。

反广群：I，§ 1，第 1 号。

与 M-集相反的， \(\mathbf{M}^0\) -集：I，§ 5，第 1 号。

反环：I，§ 8，第 3 号。

轨道：I，§ 5，第 4 号。

轨道映射：I，§ 5，第 4 号。

阶，无限阶元素：I，§ 4，第 10 号。

轮换的阶：I，§ 5，第 7 号。

形式幂级数关于某些未定元的阶：III，§ 2，第 11 号。

群的阶：I，§ 4，第 1 号。

群中元素的阶：I，§ 4，第 10 号。

方阵的阶：II，§ 10，第 7 号。

形式幂级数的阶，全序：III，§ 2，第 11 号。

有序序列：I，§ 1，第 2 号。

有序序列，相似的：I，§ 1，第 2 号。

原点：I，§ 2，第 1 号。

原点，在仿射空间中的选取：II，§ 9，第 1 号。

正交元素，正交集合：II，§ 2，第 3 号。

投影算子的正交族：II，§ 1，第 8 号。

正交子模，关于 E（相应地 E*）的子集：II，§ 2，第 4 号。

平行线性簇：II，§ 9，第 3 号。

平行四边形：II，§9，习题1。

仿射直线的方向参数：II，§ 9，第 3 号。

偏导数：III，§ 10，第 11 号。

偏分次：II，§ 11，第 1 号。

过渡，过渡矩阵：II，§ 10，第 8 号。

可置换元素：I，§ 1，第5号。

置换，偶置换，奇置换：I，§ 5，第 7 号。

平面，仿射：II，§ 9，第 1 和第 3 号。

向量空间中过 0 的平面：II，§ 7，第 3 号。

平面，射影：II，§ 9，第 5 号。

仿射空间中的点：II，§ 9，第 1 号。

射影空间的点：II，§ 9，第 5 号。

点，仿射无关：II，§ 9，第 3 号。

无穷远点：II，§ 9，第 8 号。

多项式，特征的，K-代数中元素的：III，§ 9，第 3 号。

多项式，自同态的特征多项式：III，§ 8，第 11 号。

多项式，标量关于模的特征多项式：III，§ 9，第 1 号。

不含 \(\mathbf{X}^{\nu}\) 项的多项式：III，§ 2，第 9 号。

多项式恒等式：III，§ 2，第 9 号。

\(n\) 次多项式：III，§ 2，第 9 号。

多项式关系子：III，§ 2，第 9 号。

无常数项的多项式：III，§ 2，第 9 号。

关于一族未定元的多项式，系数在环中：III，§ 2，第 9 号。

正有理整数：I，§ 2，第 5 号。

正有理数：I，§ 9，第 4 号。

线性映射的外幂：III，§ 7，第 4 号。

矩阵的外幂：III，§ 8，第 5 号。

模的外幂：III，§ 7，第 4 号。

幂，结合律下的 \(n\) 次幂：I，§ 1，第 3 号。

幂，对称的，线性映射的、模的：III，§ 6，第 3 号。

幂，张量的，线性映射的：III，§ 5，第 2 号。

幂，张量的，模的：III，§ 5，第 1 号。

群的表现：I，§ 7，第 6 号。

代数的表出：III，§ 2，第 8 号。

有限表现的（群）：I，§ 7，习题 16。

传递到商时的保持性：I，§ 1，第 6 号。

素理想：I，§ 9，第3号。

素数：I，§ 4，第 10 号。

素数，互素（整数）：I，§ 8，第 10 号。

自由群中的本原元素：I，§ 7，习题 26。

由单个元素生成的自由广群中的本原元素：I，§ 7，习题 7。

分次双代数的本原元：III，§ 11，第 8 号。

主 G-集，齐次的：I，§ 5，第 6 号。

主理想：I，§ 8，第 6 号。

主列：I，§ 4，习题 17。

G 下的主集，齐次的：I，§ 5，第 6 号。

线性延拓原理：II，§ 1，第 7 号。

问题，线性：II，§ 2，第 8 号。

积代数：III，§ 1，第 4 号。

积，矩阵的分块积：II，§ 10，第 5 号。

积，交叉积：III，§ 2，习题 11。

积，外积，一个 \(p\) 向量与一个 \(q\) 向量的外积：III，§ 7，第 1 号。

积，G 与 F 的外半直积，相对于 \(\tau\) ：I，§ 6，第 1 号。

积，纤维积，带算子群的：I，§ 4，第 8 号。

积，纤维积，模的：II，§ 1，习题 4。

积，自由的，代数的：III，§ 5，习题 6。

积，自由的，群的：I，§ 7，第 5 号。

积，分次张量积，两个分次模的：II，§ 11，第 5 号。

积，分次张量积，分次代数的：III，§ 4，第 7 和第 9 号。

积群，内积群，一族商群的：I，§ 4，第 8 号。

带算子群的积群：I，§ 4，第 8 号。

子群的内直积、直积、积：I，§ 4，第 9 号。

积广群：I，§ 1，第 1 号。

K'-结构的积：II，§ 8，第 3 号。

合成律的积：I，§ 1，第 1 号。

按映射计算的矩阵的积：II，§ 10，第 2 号。

在环中计算的矩阵的积：II，§ 10，第 3 号。

模的积：II，§ 1，第 5 号。

多重模的积：II，§ 1，第 14 号。

算子与元素的积：I，§ 3，第 1 号。

有序序列的积：I，§ 1，第 2 号。

两个元素的积：I，§ 1，第 1 号。

双边理想的积：I，§ 8，第 9 号。

积，（右、左）内乘：III，§ 11，第 6 和第 7 号。

积环：I，§ 8，第 10 号。

多重线性映射的对称积：III，§ 11，第 2 号。

积，张量积，一族 \(\mathbf{Z}\) -模关于三元组 \((c, p, q)\) 的：II，§ 3，第 9 号。

代数的张量积：III，§ 4，第 1 号。

无限族代数的张量积：III，§ 4，第 5 号。

代数基的张量积：III，§ 4，第 5 号。

余代数的张量积：III，§ 11，第 1 号。

两个基的张量积：II，§ 3，第 7 号。

两个元素的张量积：II，§ 3，第 1 号。

两个线性映射的张量积：II，§ 3，第 2 号。

交换环上两个矩阵的张量积：II，§ 10，第 10 号

两个模的张量积：II，§ 3，第 1 号。

两个多重模的张量积：II，§ 3，第 4 号。

两个半线性映射的张量积：II，§ 3，第 3 号。

投影同态：I，§ 4，第 8 号。

射影域：II，§ 9，第 9 号。

射影群：II，§ 9，第 10 号。

射影超平面、射影平面：II，§ 9，第 7 号。

射影直线：II，§ 9，第 5 号。

射影映射：II，§ 9，第 10 号。

射影模：II，§ 2，第 2 号。

射影空间：II，§ 9，第 5 号和第 11 号。

射影自由的、射影相关的、族：II，§ 9，第 7 号。

投影算子：II，§ 1，第 8 号。

伪模，左、右：II，附录，第 2 号。

伪环：I，§ 8，第 1 号。

平方为零的伪环：I，§ 8，第 3 号。

纯 \(p\) -向量：III，§ 11，第 13 号。

纯四元数：III，§ 2，习题 3。

二次代数：III，§ 2，第 3 号。

拟群：I，§ 3，习题 6。

四元数，纯：III，§ 2，习题 3。

四元数代数：III，§ 2，第 5 号。

四元数群：I，§ 6，习题 4。

商代数：III，§ 1，第 2 号。

商分次：II，§ 11，第 3 号。

带算子的商群：I，§ 4，第 4 号。

商律：I，§ 1，第 6 号。

商广群：I，§ 1，第 11 号。

商模、商向量空间：II，§ 1，第 3 号。

商多重模：II，§ 1，第 14 号。

作用的商：I，§ 3，第 3 号。

商环：I，§ 8，第 7 号。

带算子的群的合成列的商：I，§ 4，第 7 号。

自由群的秩：I，§ 7，习题 14。

向量空间的线性映射的秩：II，§ 7，第 4 号。

域上线性方程组的秩：II，§ 7，第 6 号。

域上矩阵的秩：II，§ 10，第 12 号。

仿射线性映射的秩：II，§ 9，第 4 号。

向量空间张量积中元素的秩：II，§ 7，第 8 号。

向量空间的半线性映射的秩：II，§ 7，第 4 号。

整环上模的子集的秩：II，§ 7，第 10 号。

向量空间的子集的秩：II，§ 7，第 2 号。

位似比：II，§ 1，第 1 号。

有理整数：I，§ 2，第 5 号。

有理数：I，§ 9，第 4 号。

有理数域：I，§ 9，第 4 号。

在子域上有理（线性形式）：II，§ 8，第 4 号。

在子域上有理（线性映射）：II，§ 8，第 3 号。

在子域上有理（子空间、向量）：II，§ 8，第 2 号。

融合和中元素的约化分解：I，§ 7，第 3 号。

自反模：II，§ 2，第 7 号。

正则、左正则、右正则元素：I，§ 2，第 2 号。

p-正则元素：I，§ 6，习题 28。

相关系统、相关子集：II，§ 1，第 11 号，及 § 9，第 7 号。

关系、等价关系，与合成律相容：I，§ 1，第 6 号。

关系、等价关系，与作用相容：I，§ 3，第 3 号。

关系、等价关系，由有序对族生成：I，§ 1，第 6 号。

关系、等价关系，左、右，与合成律相容：I，§ 3，第 3 号。

乘法表中的交换关系：III，§ 1，第 7 号。

格拉斯曼关系：III，§ 11，第 13 号。

线性关系模：II，§ 1，第 11 号。

互素整数：I，§ 8，第 11 号。

关系子：I，§ 7，第 6 号及 III，§ 2，第 8 号。

泛关系子：III，§ 2，第 8 号。

关系子理想：III，§ 3，第 8 号。

一个表现的诸关系子：I，§ 7，第 6 号及 III，§ 2，第 8 号。

多项式关系子：III，§ 2，第 9 号。

半群的剩余：I，§ 2，习题 11。

剩余有限群：I，§ 5，习题 5。

幺半群的限制代数：III，§ 2，第 10 号。

限制和，内限制和：I，§ 4，第 9 号。

群的限制和：I，§ 4，第 9 号。

关于子群的群的限制和：I，§ 4，第 9 号。

标量限制（由此得到的代数）：III，§ 1，第 5 号。

（将元素代入不定元而）得到的（元素）：I，§ 7，第 5 号。

扩张的收缩：I，§ 6，第 1 号。

线性方程组的右端项，线性系统的右端项：II，§ 2，第 8 号。

环：I，§ 8，第 1 号。

环，布尔：I，§ 9，习题 8。

环，交换的：I，§ 8，第 1 号。

环（分次的，正次数分次的，双分次的）：II，§ 11，第 2 号。

环，分次商环：II，§ 11，第 3 号。

通过添加单位元所得的环：II，附录，第 1 号。

交换群的自同态环：I，§ 8，第 4 号。

分母在 S 中的分式环：I，§ 8，第 12 号。

模 \(n\) 整数环：I，§ 8，第 11 号。

关于一个自同态与一个导子不可交换的多项式环：III，§ 10，习题 3。

环，反环：I，§ 8，第 3 号。

环，乘积：I，§ 8，第 10 号。

环，商环：I，§ 8，第 7 号。

环，全分式环：I，§ 8，第 2 号。

环，零环：I，§ 8，第 3 号。

矩阵的行：II，§ 10，第 1 号。

法则，符号：I，§ 8，第 1 号。

标量：II，§ 1，第 1 号。

标量线性方程：II，§ 2，第 8 号。

标量矩阵：II，§ 10，第 7 号。

施赖埃尔关于合成列的定理：I，§ 4，第 7 号。

扩张的截面：I，§ 6，第 2 号。

群的半直积：I，§ 6，第 1 号。

半群，左、右：I，§ 2，习题 11。

半同态，代数，代数 \(\rho\) -同态：III，§ 1，第 5 号。

半线性映射：II，§ 1，第 13 号。

序列，正合：II，§ 1，第 4 号。

序列，有序：I，§ 1，第 2 号。

序列，分裂正合：II，§ 1，第 9 号。

序列，相似有序：I，§ 1，第 2 号。

级数，形式幂级数的代数：III，§ 2，第 11 号。

列，合成：I，§ 4，第 6 号。

列，导出的：I，§ 6，第 4 号。

列，等价分解的：I，§ 4，第 7 号。

列，更细的合成：I，§ 4，第 7 号。

级数，形式幂：III，§ 2，第 11 号。

列，若尔当-赫尔德：I，§ 4，第 7 号。

列，下中心：I，§ 6，第 3 号。

列，正规：I，§ 4，习题 17。

列，主：I，§ 4，习题 17。

G-集，齐次的（G 为算子群）：I，§ 5，第 5 号。

G-集，主齐次的，G 下的主齐次集：I，§ 5，第 6 号。

集合，生成，群的：I，§ 4，第 3 号。

集合，生成，广群的：I，§ 1，第 4 号。

M-集（M 为算子幺半群）：I，§ 5，第 1 号。

\(\mathbf{M}^0\) -集，与 M-集相反的：I，§ 5，第 1 号。

分次群的次数集：II，§ 11，第 1 号。

满足极大条件、极小条件的子群集合：I，§ 4，习题 15。

集合，正交：II，§ 2，第 4 号。

G-集，弱等价：I，§ 5，习题 26。

移位，平移：II，§ 11，第 2 号。

置换的符号差：I，§ 5，第 7 号。

有理数的符号：I，§ 9，第 4 号。

相似自同态：III，§ 8，习题 26。

相似矩阵：II，§ 10，第 9 号。

相似有序序列：I，§ 1，第 2 号。

单群：I，§ 4，第 4 号。

单传递作用：I，§ 5，第 6 号。

除环：I，§ 9，第 1 号。

斜分次双代数：III，§ 11，第 4 号。

反对称张量：III，§ 7，第 4 号。

反对称化张量：III，§ 7，第 4 号。

分次代数的斜张量积：III，§ 4，第 7 和第 8 号。

线性方程、线性方程组的解：II，§ 2，第 8 号。

解，平凡的，齐次线性方程的零解：II，§ 2，第 8 号。

可解群，\(n\) 类可解群：I，§ 6，第 4 号。

空间，仿射，附于向量空间：II，§ 9，第 1 号。

空间，典范射影，与向量空间相伴：II，§ 9，第 8 号。

空间，射影：II，§ 9，第 5 和第 11 号。

空间，射影，由向量空间导出：II，§ 9，第 5 号。

空间，商向量空间，商空间：II，§ 1，第 3 号。

空间，右、左、向量，域上的：II，§ 1，第 1 号。

空间，向量，与整环上模相伴：II，§ 7，第 10 号。

空间，向量，通过在仿射空间中取原点得到：II，§ 9，第 1 号。

空间，向量，仿射空间的平移：II，§ 9，第 1 号。

分裂正合序列：II，§ 1，第 10 号。

稳定（算子、算子集合）一个子集：I，§ 5，第 2 号。

严格稳定（算子、算子集合）一个子集：I，§ 5，第 2 号。

带算子群的稳定子群：I，§ 4，第 3 号。

稳定子集：I，§ 1，第 4 号。

严格负、严格正有理整数：I，§ 2，第 5 号。

严格负、严格正有理数：I，§ 9，第 4 号。

严格稳定化子：I，§ 5，第 2 号。

严格稳定一个子集：I，§ 5，第 2 号。

严格传送子：I，§ 5，第 2 号。

强结合子集：III，附录，第 1 号。

结构，模（多重模）的，相容的：II，§ 1，第 14 号。

结构，射影空间：II，§ 9，第 11 号。

K'-结构，诱导的：II，§ 8，第 2 号。

向量 K-空间上的 K'-结构：II，§ 8，第 1 号。

K'-结构，积：II，§ 8，第 3 号。

子代数：III，§ 1，第 2 号。

由子集生成的子代数：III，§ 1，第 2 号。

分次子代数：III，§ 3，第 2 号。

子域：I，§ 9，第 1 号。

由子集生成的子域：I，§ 9，第 1 号。

特征子群：I，§ 5，第 3 号。

不变稳定子群，不变子群：I，§ 4，第 4 号。

正规稳定子群，正规子群：I，§ 4，第 4 号。

稳定子群，由子集生成的：I，§ 4，第 3 号。

稳定子群：I，§ 4，第 3 号。

西罗子群，西罗 p-子群：I，§ 6，第 6 号。

子广群：I，§ 1，第 4 号。

由子集生成的子广群：I，§ 1，第 4 号。

子广群，含幺的：I，§ 2，第 1 号。

子广群，含幺的，由子集生成的：I，§ 2，第 1 号。

子矩阵：II，§ 10，第 1 号。

子模：II，§ 1，第 3 号。

模直和的分量子模：II，§ 1，第 6 号。

由族生成的子模：II，§ 1，第 7 号。

分次子模：II，§ 11，第 3 号。

不可约子模：II，§ 2，习题 16。

子模，与 E（相应地 E*）的子集正交（或完全正交）的：II，§ 2，第 4 号。

整环上模的挠子模：II，§ 7，第 10 号。

零子模：II，§ 1，第 3 号。

互补子模：II，§ 1，第 9 号。

子多重模：II，§ 1，第 14 号。

子环：I，§ 8，第 5 号。

由子集生成的子环：I，§ 8，第 5 号。

分次子环：II，§ 11，第 3 号。

仿射子集：II，§ 9，第 3 号。

中心化某子集的子集：I，§ 5，第 3 号。

自由子集，相关子集：II，§ 1，第 11 号。

子集，齐次的，关于形式幂级数中某些不定元次数为 \(p\) 的：III，§ 2，第 11 号。

使一个子集正规化的子集：I，§ 5，第 3 号。

子集，稳定的：I，§ 1，第 4 号和 § 3，第 2 号。

子集，稳定的，由一个子集生成的：I，§ 1，第 4 号和 § 3，第 2 号。

子集，强结合的：III，附录，第 1 号。

子集，对称的：I，§ 4，第 1 号。

子集，共轭的：I，§ 5，第 4 号。

与外代数的一个齐次元素相伴的子空间：III，§ 7，第 2 号。

与对称代数的一个齐次元素相伴的子空间：III，§ 6，第 2 号。

与张量代数的一个齐次元素相伴的子空间：III，§ 5，第 2 号。

与向量空间的张量积的一个元素相伴的子空间：II，§ 7，第 8 号。

在一个子域上有理的子空间：II，§ 8，第 2 号。

子空间，向量子空间：II，§ 1，第 3 号。

和，融合的，模的：II，§ 1，习题 5。

和，融合的，幺半群的：II，§ 7，第 3 号。

和，直和：I，§ 4，第 9 号。

和，一个子模族的直和：II，§ 1，第 8 号。

和，一个子模族的外直和：II，§ 1，第 6 号。

和，子群的内限制和：I，§ 4，第 9 号。

和，幺半的：I，§ 7，第 3 号。

一个有限支撑元素族的和：I，§ 2，第 1 号。

一个左理想族、右理想族的和：I，§ 8，第 6 号。

一个子模族的和：II，§ 1，第 7 号。

一个有序序列的和：I，§ 1，第 2 号。

两个元素的和：I，§ 1，第 1 号。

两个矩阵的和：II，§ 10，第 2 号。

和，关于子群的群的限制和，群的限制和：I，§ 4，第 9 号。

超可解群：I，§ 6，习题 26。

互补子模：II，§ 1，第 9 号。

一个轮换的支集：I，§ 5，第 7 号。

一个族的支集：I，§ 2，第 1 号。

在矩阵中删去列、行：II，§ 10，第 1 号。

西罗子群：I，§ 6，第 6 号。

符号，克罗内克：II，§ 1，第 11 号。

模的对称代数：III，§ 6，第 1 号。

对称群：I，§ 4，第 1 号。

对称多重线性映射：III，§ 6，第 3 号。

线性映射的对称幂：III，§ 6，第 3 号。

模的对称幂：III，§ 6，第 3 号。

两个多重线性映射的对称积：III，§ 11，第 2 号。

群中的对称子集：I，§ 4，第 1 号。

对称张量：III，§ 6，第 3 号。

张量的对称化：III，§ 6，第 3 号。

系统，正向的，分次代数的：III，§ 3，第 3 号。

系统，正向的，广群的：I，§ 10，第 3 号。

系统，正向的，模的：II，§ 6，第 2 号。

系统，正向的，环（群、域）的：I，§ 10，第 3 号。

系统，自由的，相关系统：II，§ 1，第 11 号。

系统，生成，群的：I，§ 4，第 3 号。

系统，生成，广群的：I，§ 1，第 4 号。

系统，生成，模的：II，§ 1，第 7 号。

系统，生成，射影空间的：II，§ 9，第 8 号。

系统，生成，含幺广群的：I，§ 2，第 1 号。

系统，生成，代数的：III，§ 1，第 2 号。

系统，逆向的，代数的：III，§ 1，第 6 号。

系统，逆向的，广群的：II，§ 10，第 1 号。

系统，逆向的，模的：II，§ 6，第 1 号。

系统，逆向的，环（群，域）的：I，§ 10，第 1 号。

向量子空间的方程组：II，§ 7，第 5 号。

因子组：III，§ 2，习题 11。

点的齐次坐标系：II，§ 9，第 6 号。

线性方程组，线性系统，齐次线性系统：II，§ 2，第 8 号。

平凡交换因子组：III，§ 4，第 7 号。

矩阵的对角表，下三角，上三角：II，§ 10，第 7 号。

代数的乘法表：III，§ 1，第 7 号。

矩阵的方表：II，§ 10，第 5 号。

模的张量代数：III，§ 5，第 1 号。

张量，逆变张量，协变张量，混合张量：III，§ 5，第 6 号。

(I, J) 型张量：III，§ 5，第 6 号。

反对称张量：III，§ 7，第 4 号。

对称张量：III，§ 6，第 3 号。

线性映射的张量幂：III，§ 5，第 2 号。

模的张量幂：III，§ 5，第 1 号。

形式幂级数的常数项：III，§ 2，第 11 号。

多项式的常数项：III，§ 2，第 9 号。

多项式中 \(\mathbf{X}^{\alpha}\) 的项：III，§ 2，第 9 号。

形式幂级数的项：III，§ 2，第 11 号。

多项式的项：III，§ 2，第 9 号。

和的一项：I，§ 1，第 2 号。

关于形式幂级数中某些不定元次数为 \(p\) 的项：III，§ 2，第 11 号。

形式幂级数中全次数为 \(p\) 的项：III，§ 2，第 11 号。

定理，结合性：I，§ 1，第 3 号。

定理，凯莱-哈密顿：III，§ 8，第 11 号。

定理，交换性：I，§ 1，第 5 号。

定理，德萨格：II，§ 9，习题 15。

定理，埃尔德什-卡普兰斯基：II，§ 7，习题 3。

定理，射影几何的基本定理：II，§ 9，习题 16。

定理，霍尔：I，§ 6，习题 7。

定理，若尔当-赫尔德：I，§ 4，第 7 号。

定理，卡普兰斯基：II，§ 2，习题 2。

定理，克鲁尔：I，§ 8，第 6 号。

定理，尼尔森-施赖埃尔：I，§ 7，习题 20。

完全四边形的定理：II，§ 9，习题 13。

定理，帕普斯：II，§ 9，习题 14。

定理，施赖埃尔：I，§ 4，第 7 号。

挠元，挠模，挠子模：II，§ 7，第 10 号。

无挠模：II，§ 7，第 10 号。

幺半群的全代数：III，§ 2，第 10 号。

全分次：II，§ 11，第 1 号。

完全正交于一个子集的（子模）：II，§ 2，第 4 号。

迹，凯莱：III，§ 2，第 4 号。

二次代数中的迹：III，§ 2，第 3 号。

矩阵的迹：II，§ 10，第 11 号。

K-代数中元素的迹：III，§ 9，第 3 号。

自同态的迹：II，§ 4，第 3 号。

关于模的标量的迹：III，§ 9，第 1 号。

传递作用：I，§ 5，第 5 号。

传递性公式：III，§ 9，第 4 号。

仿射空间中的平移，平移空间：II，§ 9，第 1 号。

平移，左、右：I，§ 2，第 2 号。

平移，左、右（幺半群对自身的作用）：I，§ 5，第 1 号。

传送子，严格传送子：I，§ 5，第 2 号。

线性映射的转置，半线性映射的转置：II，§ 2，第 5 号。

矩阵的转置：II，§ 10，第 1 号。

对换：I，§ 5，第 7 号。

平延：II，§ 10，习题 11。

三角、下三角、上三角矩阵：II，§ 10，第 7 号。

平凡扩张：I，§ 6，第 1 号。

平凡分次：II，§ 11，第 1 号。

平凡同态：I，§ 2，第 1 号。

平凡作用：I，§ 5，第 2 号。

齐次线性方程的平凡解：II，§ 2，第 8 号。

平凡的交换因子组：III，§ 4，第 7 号。

双边理想：I，§ 8，第 6 号及 III，§ 1，第 2 号。

型，指数的，群的：I，§ 7，习题 39。

型 \(\Delta\) 的分次：II，§ 11，第 1 号。

幺模自同态：III，§ 8，第 1 号。

幺模群：III，§ 8，第 9 号。

幺模矩阵：III，§ 8，第 3 号。

\(p\) -幂单元素：I，§ 6，习题 28。

代数的单位元：III，§ 1，第 1 号。

单位元，左、右，在群胚中的：I，§ 4，习题 23。

单位，广群的单位元：I，§ 2，第 1 号。

含单位元代数：III，§ 1，第 1 号。

保单位代数同态，态射：III，§ 1，第 1 号。

保单位同态：I，§ 2，第 1 号。

含幺广群：I，§ 2，第 1 号。

矩阵单位：II，§ 10，第 3 号。

由生成系统及一族关系子所定义的泛代数：III，§ 2，第 8 号。

受恒等式约束的泛代数：III，§ 2，第 8 号。

泛关系子：III，§ 2，第 8 号。

线性方程组的未知数：II，§ 2，第 8 号。

有理数的绝对值：I，§ 9，第 4 号。

自由代数中元素的值：III，§ 2，第 8 号。

范德蒙德行列式：III，§ 8，第 6 号。

仿射线性簇，由一族生成的仿射线性簇：II，§ 9，第 3 号。

线性簇：II，§ 9，第 3、7 和 11 号。

射影线性簇，由一族生成的射影线性簇：II，§ 9，第 7 和 11 号。

平行线性簇：II，§ 9，第 3 号。

向量：II，§ 1，第 1 号。

仿射直线的方向向量：II，§ 9，第 3 号。

仿射空间的自由向量：II，§ 9，第 1 号。

\(p\) -向量：III，§ 7，第 1 号。

\(p\) -向量，纯的：III，§ 11，第 13 号。

在子域上有理的向量：II，§ 8，第 2 号。

向量空间：II，§ 1，第 1 号。

字：I，§ 7，第 2 号。

齐次元素的权：II，§ 11，第 1 号。

无不动点（自同构）：I，§ 6，习题 23。

扎森豪斯引理：I，§ 4，第 7 号。

零：I，§ 2，第 1 号。

零矩阵：II，§ 10，第 22 号。

零环：I，§ 8，第 3 号。

线性方程的零解：II，§ 2，第 8 号。

零子模：II，§ 1，第 3 号。
